\documentclass[10pt,a4paper]{amsart}
\usepackage[margin=19mm]{geometry}
\usepackage{amsmath,amssymb,mathtools}
\usepackage[scr=rsfs,cal=euler]{mathalfa}
\usepackage{xfrac}
\usepackage[unicode,hidelinks,bookmarks=false]{hyperref}
\newcommand{\ie}{i.e.\ }
\newcommand{\cf}{cf.\ }
\newcommand{\resp}{respectively\ }
\newcommand{\Subsection}{\subsection}
\providecommand{\nobreakdash}{\nobreak\hbox{-}}
\setlength{\emergencystretch}{2em}
\multlinegap0pt
\usepackage{amssymb}
\usepackage{enumerate}
\newtheorem{lemma}{Lemma}
\title{Intersecting well approximable and missing digit sets}
\author{Bing Li, Sanju Velani, Bo Wang}
\date{Source: arXiv:2512.17173v1 (CC BY 4.0)}
\begin{document}
\maketitle

\noindent\emph{Source and licence.} Intersecting well approximable and missing digit sets. Version arXiv:2512.17173v1 (CC BY 4.0), distributed by arXiv under the Creative Commons Attribution 4.0 International licence, http://creativecommons.org/licenses/by/4.0/, as recorded in the arXiv licence metadata for this exact version and shown on the abstract page https://arxiv.org/abs/2512.17173v1. Full licence notice: \texttt{licenses/arxiv-2512.17173-CC-BY-4.0.txt}.\par\smallskip

\noindent\textit{Editorial scope.} Counted unit: the article's lemma 30 (source0.tex lines 895--904). The article's complete original proof of it is delivered with it (source0.tex lines 905--947, one \texttt{proof} environment of 43 lines). No mathematics is written, repaired or re-derived: the statement and the proof are the article's own lines, in the article's own notation, and no retained line is substituted or re-flowed.  Retained uncounted as the source-local context the proof needs: lines 848--852.  Nothing was omitted from any retained span, so the package carries no omission record.  No retained line contains a citation, so the wrapper carries no bibliography and no bibliography entry.  No retained line contains a figure environment, an image-inclusion command or drawing code, so no image is delivered and the package declares no asset dependency.  Editorial formatting only, in addition: an AMS \texttt{amsart} preamble that restates the article's own declarations and macros used by the retained lines, namely the environments \texttt{lemma} on separate counters, together with the standard packages those lines need.  The retained environments print in this document as Lemma 1 in order of appearance, whereas the article prints the same items with the numbering of its own full document.  The complete native source of the article as distributed in the arXiv e-print for this exact version is bundled unchanged as \texttt{sources/191297/source0.tex}.
\begin{lemma}\label{divisible 5}  
Suppose $b$ and $t$ have the same prime divisors.  Then,  
for any $n\in\mathbb{N}$,
$$t^{n}\mid b^{\lceil \alpha_{2}n\rceil}.$$
\end{lemma}

\begin{lemma}\label{30} 
Let $b,t$ be integers with $b\geq3$ and $t\geq2$.
\begin{enumerate}
\item[(i)] Let $n,m\in\mathbb{N}$.  
Then  
$$\Gamma_{n}(\psi)\subseteq F_{n,m}(\psi).$$
\item[(ii)] Suppose $b$ and $t$ have the same prime divisors.  If $\psi:\mathbb{N}\to(0,\infty)$ satisfies $\psi(n)\leq b^{-\lceil\alpha_{2}n\rceil}$ for all $n\in\mathbb{N}$, then  
\begin{equation}\label{svweakform} \Gamma_{n}(\psi)\subseteq\left\{0\leq p\leq t^{n}:\frac{p}{t^{n}}\in E (\lceil\alpha_{2}n\rceil )\right\}.   \end{equation}
\end{enumerate}
\end{lemma}

\begin{proof}[\rm \textbf{Proof}] (i) Without loss of generality, we assume that $\Gamma_{n}(\psi)\neq\emptyset$. By definition, for every $p\in\Gamma_{n}(\psi)$, since $$C(b,D)\subseteq C_{m}(b,D)=\bigcup_{i=1}^{(\#D)^{m}}I_{m,i}  \, , $$ it follows that there exists $1\leq i\leq (\#D)^{m}$, such that 
$$ B\big(\textstyle{\frac{p}{t^{n}},\psi(n)}\big)\cap I_{m,i}\neq\emptyset.$$
We now consider two cases depending on  whether or not the endpoints $ \frac{q_{i}}{b^{m}},\frac{q_{i}+1}{b^{m}} $   of $I_{m,i}$ lie  in the ball $ B\left(\frac{p}{t^{n}},\psi(n)\right)$. 
%According to whether or not the endpoints $ \frac{q_{i}}{b^{m}},\frac{q_{i}+1}{b^{m}} $   of $I_{m,i}$ lie  in the ball $ B\left(\frac{p}{t^{n}},\psi(n)\right)$, we divide the proof into two cases.
\begin{itemize}
    \item {\em  There exists  one endpoint of $I_{m,i}$ that lies in 
$B\left(\frac{p}{t^{n}},\psi(n)\right)$.} Let us say $\frac{q_{i}}{b^{m}}\in B\left(\frac{p}{t^{n}},\psi(n)\right)$ without loss of generality. Then
$$\left|\frac{p}{t^{n}}-\frac{q_{i}}{b^{m}}\right|<\psi(n).$$
    \item {\em   Both endpoints of $I_{m,i}$ do not lie in $B\left(\frac{p}{t^{n}},\psi(n)\right)$.} In this case, we have 
$$B\big(\textstyle{\frac{p}{t^{n}},\psi(n)}\big) \subseteq I_{m,i}.$$
In particular, $\frac{p}{t^{n}}\in\left(\frac{q_{i}}{b^{m}},\frac{q_{i}+1}{b^{m}}\right)$. Thus,
$$\left|\frac{p}{t^{n}}-\frac{q_{i}}{b^{m}}\right|<\frac{1}{b^{m}}\quad {\rm and}\quad \left|\frac{p}{t^{n}}-\frac{q_{i}+1}{b^{m}}\right|<\frac{1}{b^{m}}.$$
\end{itemize}

%\noindent \text{\em Case 1}: There exists  one endpoint of $I_{m,i}$ that lies in  $B\left(\frac{p}{t^{n}},\psi(n)\right)$. Let us say $\frac{q_{i}}{b^{m}}\in B\left(\frac{p}{t^{n}},\psi(n)\right)$ without loss of generality. Then $$\left|\frac{p}{t^{n}}-\frac{q_{i}}{b^{m}}\right|<\psi(n).$$

%\noindent \text{\em Case 2}: Both endpoints of $I_{m,i}$ do not lie in $B\left(\frac{p}{t^{n}},\psi(n)\right)$. In this case, we have  $$B\Big(\frac{p}{t^{n}},\psi(n)\Big)\subseteq I_{m,i}.$$ In particular, $\frac{p}{t^{n}}\in\left(\frac{q_{i}}{b^{m}},\frac{q_{i}+1}{b^{m}}\right)$. Thus, $$\left|\frac{p}{t^{n}}-\frac{q_{i}}{b^{m}}\right|<\frac{1}{b^{m}}\quad {\rm and}\quad \left|\frac{p}{t^{n}}-\frac{q_{i}+1}{b^{m}}\right|<\frac{1}{b^{m}}.$$
\noindent On combining the above two cases, we obtain that there exists an endpoint $\frac{q}{b^{m}}\in E(m)$, such that $$\left|\frac{p}{t^{n}}-\frac{q}{b^{m}}\right|<\max\Big(\psi(n) ,\frac{1}{b^{m}}\Big).$$
Therefore, for every $p\in\Gamma_{n}(\psi)$ we have that  $p\in F_{n,m}(\psi)$ and so  this implies that 
$$\Gamma_{n}(\psi)\subseteq F_{n,m}(\psi).$$


%\bigskip

\noindent  (ii) Fix $n\in\mathbb{N}$. Without loss of generality, we suppose that $\Gamma_{n}(\psi)\neq\emptyset$. It follows from Lemma $\ref{30}$ (i), on putting $m=\lceil \alpha_{2}n\rceil$,  that  
\begin{equation}\label{contain}
\Gamma_{n}(\psi)\subseteq F_{n,\lceil \alpha_{2}n\rceil}(\psi).
\end{equation}
 Let $p\in\Gamma_{n}(\psi)$. In view of \eqref{contain}, there exists $\frac{q}{b^{\lceil \alpha_{2}n\rceil}}\in E (\lceil\alpha_{2}n\rceil)$, such that
$$\left|\frac{p}{t^{n}}-\frac{q}{b^{\lceil \alpha_{2}n\rceil}}\right|<\max\Big(\psi(n),\frac{1}{b^{\lceil \alpha_{2}n\rceil}}\Big).$$
This together with the fact that $\psi(n)\leq b^{-\lceil \alpha_{2}n\rceil}$,  implies that
$$\left|\frac{p}{t^{n}}-\frac{q}{b^{\lceil \alpha_{2}n\rceil}}\right|<\frac{1}{b^{\lceil \alpha_{2}n\rceil}}.$$
Hence,
\begin{equation}\label{inequality}
\left|q-\frac{b^{\lceil \alpha_{2}n\rceil}}{t^{n}}p\right|<1.
\end{equation}
By Lemma \ref{divisible 5}, we know that $t^{n}\mid b^{\lceil \alpha_{2}n\rceil}$ and so the left-hand side of \eqref{inequality} is an integer. The upshot of this and \eqref{inequality} is that
$$q=\frac{b^{\lceil \alpha_{2}n\rceil}}{t^{n}}p.$$ Therefore,
$\ \displaystyle{\small{\frac{p}{t^{n}}=\frac{q}{b^{\lceil \alpha_{2}n\rceil}}  } \in E(\lceil\alpha_{2}n\rceil)}  \ $
%$p/t^{n}=q/b^{\lceil \alpha_{2}n\rceil}\in E(\lceil\alpha_{2}n\rceil) $
as desired.
\textrm{}
\end{proof}

\end{document}
